% Source: https://artofproblemsolving.com/community/c3h2304256_solution_to_problem_18_of_the_amc_8
% Stable identity: aops-topic-2304256-q01
\textbf{(a)} In an arithmetic sequence $u,u+d,u+2d,u+3d,u+4d$ with five terms, the middle term is the average of the first and the last: $u+2d=\dfrac{u+(u+4d)}{2}$.

So the middle entry of the top row is $\dfrac{1+25}{2}=13$ and the middle entry of the bottom row is $\dfrac{17+81}{2}=49$. These are the first and last entries of the middle column, whose middle entry is $X$, so
\[X=\dfrac{13+49}{2}=31.\]
The answer is $\boxed{\textbf{(B) }31}$.

\textbf{(b)} The computation in (a) shows that $X=\dfrac{1+25+17+81}{4}$, the average of the four corner entries. The student's expression is also $\dfrac{1+81+17+25}{4}$, the average of the same four numbers, so it gives the same value $31$. But each diagonal on its own does not give $X$: $\dfrac{1+81}{2}=41$ and $\dfrac{17+25}{2}=21$. Indeed, filling in the array (each row from its two ends, then checking the columns) gives
\[\begin{array}{ccccc}
1&7&13&19&25\\
5&13.5&22&30.5&39\\
9&20&31&42&53\\
13&26.5&40&53.5&67\\
17&33&49&65&81
\end{array}\]
Every row and column is arithmetic (for example the second column $7,13.5,20,26.5,33$ has difference $6.5$), and the centre is $31$. The diagonal $1$, $13.5$, $31$, $53.5$, $81$ has differences $12.5$, $17.5$, $22.5$, $27.5$, and the diagonal $25$, $30.5$, $31$, $26.5$, $17$ goes up and then down, so neither diagonal is an arithmetic sequence. The student's formula is correct only because the two errors cancel: the four corners together average to $X$.
